% ============================================================
% ONE-PAGE ACADEMIC POSTER
% Topic : Secure Hardware Chips – Architecture, Attacks & Future
% Engine: pdfLaTeX  |  Compile on Overleaf (or locally)
% ============================================================
\documentclass[final]{beamer}

% ---------- poster geometry ----------
\usepackage[orientation=portrait,size=a1,scale=1.1]{beamerposter}

% ---------- general packages ----------
\usepackage[utf8]{inputenc}
\usepackage{lmodern}
\usepackage{amsmath,amssymb}
\usepackage{graphicx}
\usepackage{xcolor}
\usepackage{booktabs}       % professional table rules
\usepackage{array}          % extended column definitions
\usepackage{tikz}
\usetikzlibrary{positioning,arrows.meta,shapes.geometric,backgrounds,fit}
\usepackage{pgfplots}
\pgfplotsset{compat=1.18}
\usepackage{enumitem}
\usepackage{microtype}

% ============================================================
% COLOUR PALETTE
% ============================================================
\definecolor{PrimaryBlue}{HTML}{0A3161}
\definecolor{AccentTeal}{HTML}{1B7F8E}
\definecolor{LightGray}{HTML}{F4F6F8}
\definecolor{DarkGray}{HTML}{2C2C2C}
\definecolor{AlertRed}{HTML}{C0392B}
\definecolor{GreenOK}{HTML}{1A7A4A}

% ============================================================
% BEAMER THEME OVERRIDES
% ============================================================
\usetheme{default}
\setbeamercolor{background canvas}{bg=LightGray}
\setbeamercolor{headline}{bg=PrimaryBlue,fg=white}
\setbeamercolor{block title}{bg=AccentTeal,fg=white}
\setbeamercolor{block body}{bg=white,fg=DarkGray}
\setbeamercolor{block title alerted}{bg=AlertRed,fg=white}
\setbeamercolor{block body alerted}{bg=white,fg=DarkGray}
\setbeamercolor{block title example}{bg=GreenOK,fg=white}
\setbeamercolor{block body example}{bg=white,fg=DarkGray}
\setbeamerfont{block title}{size=\large,series=\bfseries}
\setbeamerfont{headline title}{size=\huge,series=\bfseries}
\setbeamerfont{headline author}{size=\large}
\setbeamerfont{headline institute}{size=\normalsize}

% remove navigation symbols
\setbeamertemplate{navigation symbols}{}

% ---- headline (title bar) ----
\setbeamertemplate{headline}{
  \begin{beamercolorbox}[wd=\paperwidth,ht=8cm,dp=0.5cm,
                         leftskip=1.5cm,rightskip=1.5cm]{headline}
    \vskip1.2cm
    {\usebeamerfont{headline title}%
     \color{white}Secure Hardware Chips:\\[0.3ex]
     \large Architecture, Threat Landscape \& Future Frontiers}\\[0.8em]
    {\usebeamerfont{headline author}\color{white!85!black}%
     Yashraj Singh}\\[0.3em]
    {\usebeamerfont{headline institute}\color{white!70!black}%
     Department of Electronics and Communication Engineering (ECE),
     Indian Institute of Information Technology, Allahabad\quad
     \texttt{iec2023065@iiita.ac.in}}
    \vskip0.8cm
  \end{beamercolorbox}
}

% ---- footline ----
\setbeamertemplate{footline}{
  \begin{beamercolorbox}[wd=\paperwidth,ht=1.8cm,dp=0.4cm,
                         leftskip=1.5cm,rightskip=1.5cm]{headline}
    \color{white!80!black}\small
    \textbf{Secure Hardware Chips} \hfill
    IIIT Allahabad — ECE \hfill
    \today
  \end{beamercolorbox}
}

% block style
\setbeamertemplate{block begin}{
  \begin{beamercolorbox}[rounded=true,shadow=false,leftskip=0.5em,
                         colsep*=0.5ex,sep=0.5ex]{block title}%
    \usebeamerfont{block title}\insertblocktitle
  \end{beamercolorbox}%
  \begin{beamercolorbox}[rounded=true,shadow=false,leftskip=0.5em,
                         rightskip=0.5em,colsep*=0.5ex,sep=0.5ex,
                         vmode]{block body}%
    \usebeamerfont{block body}\vskip0.3ex
}
\setbeamertemplate{block end}{\vskip0.5ex\end{beamercolorbox}\vskip0.4cm}

\setbeamertemplate{block alerted begin}{
  \begin{beamercolorbox}[rounded=true,shadow=false,leftskip=0.5em,
                         colsep*=0.5ex,sep=0.5ex]{block title alerted}%
    \usebeamerfont{block title}\insertblocktitle
  \end{beamercolorbox}%
  \begin{beamercolorbox}[rounded=true,shadow=false,leftskip=0.5em,
                         rightskip=0.5em,colsep*=0.5ex,sep=0.5ex,
                         vmode]{block body alerted}%
    \usebeamerfont{block body}\vskip0.3ex
}
\setbeamertemplate{block alerted end}{\vskip0.5ex\end{beamercolorbox}\vskip0.4cm}

\setbeamertemplate{block example begin}{
  \begin{beamercolorbox}[rounded=true,shadow=false,leftskip=0.5em,
                         colsep*=0.5ex,sep=0.5ex]{block title example}%
    \usebeamerfont{block title}\insertblocktitle
  \end{beamercolorbox}%
  \begin{beamercolorbox}[rounded=true,shadow=false,leftskip=0.5em,
                         rightskip=0.5em,colsep*=0.5ex,sep=0.5ex,
                         vmode]{block body example}%
    \usebeamerfont{block body}\vskip0.3ex
}
\setbeamertemplate{block example end}{\vskip0.5ex\end{beamercolorbox}\vskip0.4cm}

% ============================================================
\begin{document}
\begin{frame}[t]

% ============================================================
% MAIN CONTENT — two columns
% ============================================================
\begin{columns}[t,totalwidth=\linewidth]

%% ─── COLUMN 1 ────────────────────────────────────────────
\begin{column}{0.49\linewidth}

  % ---- Why Secure Hardware? ----
  \begin{block}{1 \; Why Secure Hardware?}
    Software-only defences (encryption layers, hardened OSes) remain
    vulnerable to \textbf{low-level physical and logical attacks} that
    operate below the software abstraction boundary.

    \medskip
    Secure hardware chips embed \emph{security primitives directly into
    silicon}, creating tamper-resistant environments that:
    \begin{itemize}[leftmargin=1em,itemsep=0.2ex]
      \item Isolate cryptographic keys from untrusted software
      \item Verify firmware integrity before execution
      \item Protect biometric and credential data at rest
      \item Resist physical decapsulation and probing
    \end{itemize}
  \end{block}

  % ---- Core Architectural Components ----
  \begin{block}{2 \; Core Architectural Components}
    \textbf{\color{AccentTeal}Hardware Root of Trust (RoT)}\\
    The foundational trusted element that enables secure boot, verifies
    firmware signatures, and prevents unauthorised software execution.
    All other trust in the system chains back to the RoT.

    \medskip
    \textbf{\color{AccentTeal}Cryptographic Accelerators}\\
    Dedicated processing units for AES, RSA, ECC, and SHA that improve
    throughput, reduce power, and isolate key material from software
    vulnerabilities in the main processor.

    \medskip
    \textbf{\color{AccentTeal}Secure Memory \& OTP Storage}\\
    Hardware-isolated regions and One-Time Programmable (OTP) fuses
    that safeguard keys and biometric templates against untrusted
    software access and cold-boot attacks.

    \medskip
    \textbf{\color{AccentTeal}Trusted Execution Environments (TEE)}\\
    Isolated processing spaces that partition system resources into
    \emph{Secure} and \emph{Normal} worlds, protecting sensitive
    applications (payment, biometrics, DRM) from the commodity OS.
  \end{block}

  % ---- Architecture Diagram ----
  \begin{block}{System Trust Hierarchy}
    \centering
    \begin{tikzpicture}[
        node distance=0.55cm,
        box/.style={draw, rounded corners=4pt, minimum width=9.5cm,
                    minimum height=1.0cm, font=\small\bfseries,
                    align=center},
        arrow/.style={-{Stealth[length=4pt]}, thick, color=AccentTeal}
      ]
      \node[box, fill=PrimaryBlue!90, text=white]  (rot)  {Hardware Root of Trust (RoT)};
      \node[box, fill=AccentTeal!80, text=white, below=of rot] (tee)  {Trusted Execution Environment (TEE)};
      \node[box, fill=AccentTeal!60, text=white, below=of tee] (crypt){Cryptographic Accelerators};
      \node[box, fill=AccentTeal!40, text=DarkGray, below=of crypt] (smem){Secure Memory / OTP Storage};
      \node[box, fill=LightGray,     text=DarkGray, below=of smem] (apps){Normal World Applications};
      \draw[arrow] (rot)  -- (tee);
      \draw[arrow] (tee)  -- (crypt);
      \draw[arrow] (crypt)-- (smem);
      \draw[arrow] (smem) -- (apps);
    \end{tikzpicture}
  \end{block}

  % ---- Key Technologies ----
  \begin{block}{3 \; Key Secure Chip Technologies}
    \small
    \renewcommand{\arraystretch}{1.25}
    \begin{tabular}{>{\bfseries\color{PrimaryBlue}}p{4.0cm}
                    p{12.0cm}}
      \toprule
      Technology & Architecture \& Function \\
      \midrule
      TPM
        & A discrete security processor that stores keys, measures
          boot sequences via Platform Configuration Registers (PCRs),
          and enables remote attestation. \\[0.3ex]
      ARM TrustZone
        & Processor-level architecture dividing execution into
          hardware-isolated \emph{Secure} and \emph{Normal} worlds
          with strict memory partitioning. \\[0.3ex]
      Intel SGX
        & Instruction-set extensions creating \emph{secure enclaves}—
          encrypted memory regions inaccessible even to the OS or
          hypervisor. \\[0.3ex]
      Apple Secure Enclave
        & A physically isolated coprocessor with its own OS,
          dedicated to biometric data (Face/Touch ID) and
          device encryption key management. \\[0.3ex]
      PUFs
        & Leverages microscopic manufacturing variations to generate
          unique, unclonable device fingerprints for lightweight
          hardware authentication. \\
      \bottomrule
    \end{tabular}
  \end{block}

\end{column}

%% ─── COLUMN 2 ────────────────────────────────────────────
\begin{column}{0.49\linewidth}

  % ---- Threat Landscape ----
  \begin{alertblock}{4 \; Threat Landscape \& Countermeasures}
    \small
    \renewcommand{\arraystretch}{1.2}
    \begin{tabular}{>{\bfseries\color{AlertRed}}p{3.2cm}
                    p{5.5cm}  p{6.0cm}}
      \toprule
      Attack & Mechanism & Countermeasure \\
      \midrule
      Side-Channel
        & Extracts keys via power, EM, or timing analysis
        & Data masking, constant-time execution, operation
          randomisation \\[0.3ex]
      Fault Injection
        & Voltage/clock glitching or laser pulses to cause
          computational errors
        & Redundant computations, error-detection codes, env.\
          monitors \\[0.3ex]
      Physical Probing
        & Chip decapsulation and internal structure scanning
        & Active metal shielding, tamper packaging, auto key
          erasure \\[0.3ex]
      Hardware Trojans
        & Malicious circuitry inserted during manufacturing
        & Logic testing, real-time runtime monitoring \\
      \bottomrule
    \end{tabular}
  \end{alertblock}

  % ---- Attack surface diagram ----
  \begin{block}{Attack Surface Overview}
    \centering
    \begin{tikzpicture}[
        node distance=0.45cm and 0.6cm,
        atk/.style={draw=AlertRed,  fill=AlertRed!15,  rounded corners=3pt,
                    minimum width=3.6cm, minimum height=0.75cm,
                    font=\small, align=center},
        cnt/.style={draw=GreenOK,   fill=GreenOK!15,   rounded corners=3pt,
                    minimum width=3.8cm, minimum height=0.75cm,
                    font=\small, align=center},
        chip/.style={draw=PrimaryBlue, fill=PrimaryBlue!20, rounded corners=5pt,
                     minimum width=3.6cm, minimum height=3.8cm,
                     font=\small\bfseries, align=center},
        arr/.style={-{Stealth[length=4pt]}, thick}
      ]
      \node[chip] (chip) {Secure\\Hardware\\Chip};
      \node[atk, left=2.0cm of chip, yshift=1.2cm]  (a1) {Side-Channel};
      \node[atk, left=2.0cm of chip, yshift=0.4cm]  (a2) {Fault Injection};
      \node[atk, left=2.0cm of chip, yshift=-0.4cm] (a3) {Probing};
      \node[atk, left=2.0cm of chip, yshift=-1.2cm] (a4) {HW Trojans};
      \node[cnt, right=2.0cm of chip, yshift=1.2cm]  (c1) {Const-time Exec.};
      \node[cnt, right=2.0cm of chip, yshift=0.4cm]  (c2) {Redundancy};
      \node[cnt, right=2.0cm of chip, yshift=-0.4cm] (c3) {Active Shielding};
      \node[cnt, right=2.0cm of chip, yshift=-1.2cm] (c4) {Logic Testing};
      \foreach \a/\c in {a1/c1, a2/c2, a3/c3, a4/c4}{
        \draw[arr, color=AlertRed]  (\a.east)  -- (chip.west|-\a.east);
        \draw[arr, color=GreenOK]   (chip.east|-\c.west) -- (\c.west);
      }
    \end{tikzpicture}
  \end{block}

  % ---- Future Frontiers ----
  \begin{exampleblock}{5 \; Future Frontiers}
    \textbf{\color{GreenOK}Post-Quantum Cryptography (PQC)}\\
    Hardware accelerators for NIST PQC standards (CRYSTALS-Kyber,
    CRYSTALS-Dilithium) are being embedded in secure chips to maintain
    performance parity against quantum threats.

    \medskip
    \textbf{\color{GreenOK}AI-Driven Anomaly Monitoring}\\
    Embedded ML models analyse real-time on-chip telemetry (power
    profiles, bus activity) to detect subtle side-channel or
    fault-injection campaigns before data is compromised.

    \medskip
    \textbf{\color{GreenOK}Secure Chiplets \& Self-Healing Hardware}\\
    Multi-die chiplet assemblies require secured inter-chiplet fabrics
    (e.g., UCIe). Adaptive circuits introduce deliberate jitter when
    attacks are detected, preventing leakage without system shutdown.

    \medskip
    \textbf{\color{GreenOK}Blockchain-Based Device Identity}\\
    PUF-derived hardware fingerprints registered on distributed ledgers
    enable scalable IoT authentication without centralised PKI.
  \end{exampleblock}

  % ---- Key Takeaways ----
  \begin{block}{6 \; Key Takeaways}
    \begin{itemize}[leftmargin=1.2em, itemsep=0.4ex]
      \item[\color{PrimaryBlue}\textbullet]
        Software security alone is insufficient against physical
        adversaries; \textbf{hardware roots of trust} are essential.
      \item[\color{PrimaryBlue}\textbullet]
        TEEs (TrustZone, SGX, Secure Enclave) form the practical
        foundation for protecting high-value secrets in consumer and
        enterprise devices.
      \item[\color{PrimaryBlue}\textbullet]
        Side-channel and fault-injection attacks exploit the physics
        of computation; countermeasures must be co-designed with silicon.
      \item[\color{PrimaryBlue}\textbullet]
        PUFs offer a compelling path to \emph{unclonable} device
        identity rooted in manufacturing randomness.
      \item[\color{PrimaryBlue}\textbullet]
        The convergence of PQC, AI monitoring, and chiplet
        architectures defines the next decade of secure hardware research.
    \end{itemize}
  \end{block}

  % ---- PQC Complexity visual ----
  \begin{block}{Algorithm Complexity Comparison}
    \centering\small
    \begin{tikzpicture}
      \begin{axis}[
          ybar,
          width=13cm, height=6cm,
          bar width=0.6cm,
          ylabel={Relative HW Cost (a.u.)},
          ylabel style={font=\scriptsize},
          symbolic x coords={AES-256, RSA-2048, ECC-256,
                             Kyber-768, Dilithium3},
          xtick=data,
          x tick label style={font=\scriptsize, rotate=30, anchor=east},
          ymin=0, ymax=9,
          ytick={0,2,4,6,8},
          yticklabel style={font=\scriptsize},
          grid=major,
          grid style={dotted,gray!50},
          enlarge x limits=0.12,
          axis background/.style={fill=white},
          every axis plot/.style={draw=none}
        ]
        \addplot[fill=AccentTeal!80]
          coordinates {(AES-256,1) (RSA-2048,5) (ECC-256,2.5)
                       (Kyber-768,3.5) (Dilithium3,7)};
      \end{axis}
    \end{tikzpicture}\\[0.3ex]
    {\scriptsize\color{DarkGray}%
     Post-quantum algorithms (Kyber, Dilithium) require
     dedicated accelerators to stay within power budgets.}
  \end{block}

\end{column}

\end{columns}
\end{frame}
\end{document}
